\section{Conclusion and Recommendations}

As \textit{Lakad} was fully developed and the evaluation process was completed, this chapter presents a summary of the key findings obtained, and conclusions were drawn based on the results of the system evaluation.

\subsection{Conclusion}

The developed system successfully solves both trip planning problems and the inconsistent promotion of tourist spots in the Province of Bulacan. This study proved that data-driven tools improve tourism accessibility for local areas when using mathematical optimization models together with modern mobile technologies. 

The Traveling Salesman Problem (TSP) algorithms comparison showed that Simulated Annealing (SA) works best for mobile implementation. The mobile environment could use exact methods and complex genetic algorithms to achieve higher precision but their computational demands make them impractical for real-time use. A real-time mobile environment needs fast solutions and SA provides this with high solution quality and fast execution. The Open Source Routing Machine (OSRM) serves as an essential component for the system scalability because it allows the system to execute matrix calculation for personalized itinerary generation without any limitations.

Implementing Adaptive Genetic Algorithm with Dynamic Mutation and Crossover Probabilities (AGAM) makes the personalized itinerary generation become possible through balancing multiple objectives, such as user interest, ratings, and travel distance. \textit{Lakad} generates a personalized itinerary that matches the user preferences while taking into consideration their specific constraints.

Finally, \textit{Lakad} received an excellent rating in system evaluation from both tourist and IT professionals because it functions as a technical system and as a tool that benefits society. The mobile application serves as a tool for the Provincial History, Arts, Culture, and Tourism Office (PHACTO)  to promote both popular and lesser-known attractions in the Bulacan. The study shows that \textit{Lakad} delivers an effective user-focused solution for sustainable tourism development  in Bulacan though its algorithms and localized data integration.

\subsection{Recommendations}
To further enhance \textit{Lakad}, future development should expand its database to include more destinations, accommodations, and pasalubong centers, while scaling its geographic reach to other provinces and developing an iOS version for wider accessibility. The system’s capabilities should be improved to support multi-day tour planning, integrate real-time traffic data with public transportation APIs, and immersive virtual exploration features. Future researchers should explore hybrid metaheuristic algorithms or different approach to potentially yield a much better routing and optimization. Lastly, migrating the system’s infrastructure to a self-hosted server rather than relying on external database and routing APIs.

